%%%%%%%%%%%%%%%%%%%%%%%%%%%%%%%%%%%%%%%%
% Lokalna verzija diplomske naloge v samostojnem LaTeX repozitoriju.
% Datoteka temelji na uradni FRI LaTeX predlogi.
%
% !BIB program = biber
%%%%%%%%%%%%%%%%%%%%%%%%%%%%%%%%%%%%%%%%
\DocumentMetadata{
  pdfstandard = A-4,
  pdfversion  = 2.0,
  lang        = sl-SI,
  colorprofiles = {sRGB_v4_ICC_preference.icc}
}

\documentclass[a4paper,12pt,openright]{book}
%\documentclass[a4paper,12pt,openright,draft]{book}

%%%%%%%%%%%%%%%%%%%%%%%%%%%%%%%%%%%%%%%%
%  PODATKI O DIPLOMI
%%%%%%%%%%%%%%%%%%%%%%%%%%%%%%%%%%%%%%%%

\newcommand{\tauthor}{Tomi Božak}
\newcommand{\temail}{tb85088@student.uni-lj.si}
\newcommand{\tmentor}{izr. prof. dr. Lovro Šubelj}
\newcommand{\tsomentor}{doc. dr. Gašper Slapničar}

\newcommand{\tstopnja}{UNIVERZITETNI ŠTUDIJSKI PROGRAM\\ PRVE STOPNJE\\ RAČUNALNIŠTVO IN MATEMATIKA}

\newcommand{\tinstFRI}{Fakulteta za računalništvo in informatiko}
\newcommand{\tinstFMF}{Fakulteta za matematiko in fiziko}

\newcommand{\tvrsta}{Diplomska naloga na univerzitetnem programu prve stopnje}

% TODO: opis teme prepiši iz STUDIS-a, ko bo končna verzija potrjena.
\newcommand{\tdescription}{%
Podatki, pridobljeni s sledilniki pogleda, se pogosto uporabljajo za klasifikacijo čustvenih stanj posameznikov med izvajanjem različnih nalog. V diplomskem delu se osredotočite na problem napovedovanja pozitivne oziroma negativne valence. V ta namen predlagajte ustrezno arhitekturo grafovskih nevronskih mrež ter uspešnost naučenega modela primerjajte z drugimi pristopi, kot so klasični modeli strojnega učenja za tabelarične podatke ter prednaučeni kodirniki signalov. Na koncu kritično ovrednotite uspešnost pristopa, ki temelji na grafovskem modeliranju, in predlagajte smernice za prihodnje delo.
}
% \newcommand{\tdescriptionEn}{%
% TODO: Prevedi uradni opis teme diplomskega dela iz STUDIS-a. 
% }

\newcommand{\ttitle}{Grafovske nevronske mreže za klasifikacijo valence iz podatkov sledilnika pogleda}
\newcommand{\ttitleEn}{Graph Neural Networks for Valence Classification from Eye-Tracking Data}

\newcommand{\tkeywords}{grafovske nevronske mreže, časovno-prostorski grafi, podatki sledilnika pogleda, klasifikacija valence}
\newcommand{\tkeywordsEn}{graph neural networks, spatio-temporal graphs, eye-tracking data, valence classification}

\newcommand{\tzahvala}{
  Iskreno se želim zahvaliti obema mentorjema, izr. prof. dr. Lovru Šublju za usmerjanje na področju grafovskih nevronskih mrež in doc. dr. Gašperju Slapničarju za nesebično svetovanje --~nemalokrat tudi izven delovnega časa~--, dvoletno mentoriranje na raziskovalnem področju in usmerjanje tako pri večjih kot tudi manjših odločitvah.
  Odseku E9 na IJS se zahvaljujem za neprecenljive izkušnje na področju raziskovanja in podatkovnih znanosti, za prostor, računske vire in možnost, da sem diplomsko delo opravljal v sklopu raziskovalnega dela.
  Hvala sodelavcem skupine AMI za strokovne razprave ter koristne nasvete in ideje ob raziskovalnih izzivih.
  Hvala FRI in FMF za kvalitetne temelje na področju računalništva in matematike.
  Nenazadnje pa gre zahvala od srca tudi družini in partnerki Lari za večno podporo v vseh okoliščinah. 
  Brez vaše pomoči bi bil vsak korak te poti težji.
}

\newcommand{\tposvetilo}{Komanda ni uporabljena, namenoma je odkomentirana, ker posvetila ne bom imel.}


\input{podporno/preambula.tex}

\begin{document}

%%%%%%%%%%%%%%%%%%%%%%%%%%%%%%%%%%%%%%%%
% POVZETEK IN SEZNAM KRATIC
%%%%%%%%%%%%%%%%%%%%%%%%%%%%%%%%%%%%%%%%

\newcommand{\tpovzetek}{Podatki sledilnika pogleda vsebujejo časovne in prostorske odnose med meritvami. 
Časovno-prostorska grafovska predstavitev te odnose eksplicitno kodira s povezavami med vozlišči, medtem ko jih klasična tabelarična predstavitev običajno ne.
V diplomski nalogi smo preverili, ali lahko grafovske nevronske mreže te odnose uporabijo pri binarni klasifikaciji valence.
Časovna okna meritev sledilnika pogleda smo predstavili kot časovno-prostorske grafe, v katerih vozlišča ustrezajo meritvam, povezave pa opisujejo časovno bližino, prostorsko bližino in pripadnost isti fiksaciji.
Grafovske modele smo primerjali z izhodiščnima klasifikatorjema, klasičnimi modeli strojnega učenja ter prednaučenima kodirnikoma signalov.
Modele smo ovrednotili s 7-kratnim prečnim preverjanjem po subjektih na štirih množicah signalov: \emph{samo pogled}, \emph{samo zenici}, \emph{pogled in zenici} ter \emph{vsi signali}.
Rezultati kažejo, da so grafovske predstavitve informativne, saj grafovski modeli pri vseh množicah signalov presežejo naključni in večinski klasifikator in so po rezultatih primerljivi z uveljavljenimi modeli.
Najboljši grafovski model doseže $65{,}5~\%$ točnost in $64{,}4~\%$ makro F1 pri množici \emph{pogled in zenici}, vendar ostane pod najboljšim rezultatom modela \ModelName{SVM}, ki doseže $70{,}5~\%$ točnost in $69{,}7~\%$ makro F1.
V izbranih eksperimentih so najuspešnejši klasični modeli z ročno oblikovanimi značilkami, grafovsko modeliranje pa ostaja smiselna raziskovalna smer za podatke sledilnika pogleda.
}
\newcommand{\tabstract}{% če kaj spreminjaš, se drži British English!
% na koncu preveri, da je konsistenten s slovensko verzijo povzetka!
Eye-tracking data contain temporal and spatial relations between measurements that are not modelled directly by a classical tabular representation.
This thesis investigated whether graph neural networks can use these relations for binary valence classification.
Temporal windows of eye-tracking data were represented as spatio-temporal graphs, where nodes correspond to individual measurements and edges encode temporal proximity, spatial proximity, and shared fixation membership.
Graph-based models were compared with baseline classifiers, classical machine learning models, and frozen pretrained signal encoders.
The experiments used 7-fold cross-validation by subject and evaluated four signal sets: \emph{gaze}, \emph{pupil only}, \emph{gaze and pupil}, and \emph{all signals}.
The results show that graph representations are informative, since graph-based models outperform random and majority classifiers for all signal sets.
The best graph-based model achieves $65.5~\%$ accuracy and $64.4~\%$ macro F1 on the \emph{gaze and pupil} signal set, but remains below the best \ModelName{SVM} result, which reaches $70.5~\%$ accuracy and $69.7~\%$ macro F1.
In the evaluated experimental setting, classical models with hand-crafted features are therefore the strongest approach, while graph-based modelling remains a meaningful research direction for eye-tracking data.
}
\newcommand{\tseznamkratic}{{\footnotesize
\centering
\begin{tabular}{>{\raggedright\arraybackslash}p{0.18\textwidth}|>{\raggedright\arraybackslash}p{.33\textwidth}|>{\raggedright\arraybackslash}p{.33\textwidth}}
  {\bf kratica} & {\bf angleško} & {\bf slovensko} \\ \hline
  \midrule
  {\bf AUC} & area under the curve & ploščina pod krivuljo \\
  \midrule
  {\bf CUDA} & compute unified device architecture & arhitektura CUDA za računanje na grafičnih procesorjih \\
  \midrule
  {\bf EEG} & electroencephalography & elektroencefalografija \\
  \midrule
  {\bf F1} & F1 score & mera F1 \\
  \midrule
  {\bf GAT} & graph attention network & grafovska pozornostna mreža \\
  \midrule
  {\bf GCN} & graph convolutional network & grafovska konvolucijska mreža \\
  \midrule
  {\bf GELU} & Gaussian error linear unit & aktivacijska funkcija GELU \\
  \midrule
  {\bf GIN} & graph isomorphism network & grafovska izomorfna mreža \\
  \midrule
  {\bf GNN} & graph neural network & grafovska nevronska mreža \\
  \midrule
  {\bf LightGBM} & light gradient boosting machine & metoda LightGBM z gradientno spodbujenimi odločitvenimi drevesi \\
  \midrule
  {\bf LOSO} & leave-one-subject-out cross validation &  prečno preverjanje z izločitvijo enega udeleženca\\
  \midrule
  {\bf MAHNOB-HCI} & MAHNOB-HCI database & podatkovna zbirka MAHNOB-HCI \\
  \midrule
  {\bf MLP} & multilayer perceptron & večslojni perceptron \\
  \midrule
  {\bf PyG} & PyTorch Geometric & knjižnica PyTorch Geometric \\
  \midrule
  {\bf RBF} & radial basis function & radialna bazna funkcija \\
  \midrule
  {\bf ST-GCN} & spatio-temporal graph convolutional network & časovno-prostorska grafovska konvolucijska mreža \\
  \midrule
  {\bf SVM} & support vector machine & metoda podpornih vektorjev \\
\end{tabular}
}
}

\input{podporno/naslovnestrani.tex}
\mainmatter
\setcounter{page}{1}
\pagestyle{fancy}

%%%%%%%%%%%%%%%%%%%%%%%%%%%%%%%%%%%%%%%%
% GLAVNA VSEBINA
%%%%%%%%%%%%%%%%%%%%%%%%%%%%%%%%%%%%%%%%

\chapter{Uvod}
\label{sec:uvod}

Za učenje iz podatkov, kjer niso pomembne zgolj lastnosti posameznih elementov, ampak tudi odnosi med njimi, lahko uporabimo grafovske nevronske mreže (GNN)~\cite{gilmer2017mpnn,kipf2017gcn}.
Primer takšnih podatkov so podatki sledilnika pogleda.
Meritvam sledilnika pogleda lahko smiselno določimo časovne in prostorske odnose in jim s tem priredimo časovno-prostorski graf. 
Na tem grafu se z grafovsko nevronsko mrežo učimo predstavitev vozlišč in grafov, ki jih uporabimo za klasifikacijo grafa.
Podatki, uporabljeni v tej nalogi, so označeni s čustvenimi oznakami valence, zato GNN učimo binarno klasifikacijo valence, rezultate pa primerjamo z rezultati uveljavljenih negrafovskih pristopov. 
Eksperimentalno primerjavo izvedemo na štirih signalnih naborih: samo koordinatah pogleda (pogled), samo velikostih zenic (zenici), združenih koordinatah pogleda in velikostih zenic (pogled+zenici) ter celotnem naboru razpoložljivih signalov (vsi signali).

\section{Raziskovalni problem}
\label{subsec:raziskovalni-problem}

Raziskovalni problem naloge je zasnova, implementacija in ovrednotenje grafovske predstavitve podatkov sledilnika pogleda za klasifikacijo valence.


\subsection{Raziskovalna vprašanja}
\label{subsec:raziskovalna-vprasanja}

Glavno raziskovalno vprašanje je:

\begin{enumerate}[label=\textbf{RV\arabic*:}]
    \item Kakšna je klasifikacijska uspešnost grafovske nevronske mreže za binarno klasifikacijo valence iz podatkov sledilnika pogleda v primerjavi s klasifikacijsko uspešnostjo uveljavljenih negrafovskih pristopov?
\end{enumerate}

Dodatna podvprašanja so:

\begin{enumerate}[label=\textbf{PV\arabic*:}]
    \item Kako podatke sledilnika pogleda smiselno predstaviti kot časovno-prostorski graf?
    \item Kako se uspešnost klasifikatorjev razlikuje med posameznimi nabori signalov sledilnika pogleda?
    \item Kako se klasifikacijska uspešnost spreminja z večjo kompleksnostjo grafovske nevronske mreže?
\end{enumerate}

\section{Doprinosi naloge}
\label{subsec:prispevki}

Glavni doprinosi diplomske naloge so:

\begin{enumerate}
    \item zasnova časovno-prostorske grafovske predstavitve časovnih oken podatkov sledilnika pogleda;
    \item ovrednotenje arhitekturne lestvice grafovskih modelov za klasifikacijo čustvenih oznak iz podatkov sledilnika pogleda;
    \item primerjava klasifikacijske uspešnosti pri različnih signalnih naborih;
    \item primerjava grafovskih modelov z negrafovskimi osnovnimi modeli in zamrznjenimi prednaučenimi predstavitvenimi modeli.
\end{enumerate} 

\END_NOTE{Ali rabimo tele doprinose ali raje odstranimo, glede na to da se samo ponovijo raziskovalna vprašanja?}

\section{Struktura naloge}
\label{subsec:struktura}

V poglavju~\ref{sec:ozadje} predstavimo teoretično ozadje grafov, grafovskih nevronskih mrež, podatkov sledilnika pogleda in afektivnega računalništva.
V poglavju~\ref{sec:sorodna-dela} predstavimo sorodna dela, vključno z uporabo podatkovne zbirke \MAHNOB{} in metodami učenja predstavitev iz podatkov pogleda.
Poglavje~\ref{sec:podatki} opisuje uporabljene podatke, ciljne oznake in predobdelavo.
V poglavju~\ref{sec:grafovska-predstavitev} opišemo konstrukcijo časovno-prostorskega grafa.
Poglavje~\ref{sec:modeli} predstavi uporabljene modele.
V poglavju~\ref{sec:eksperimenti} opišemo eksperimentalni protokol, v poglavju~\ref{sec:rezultati} pa rezultate glavne primerjave modelov in signalnih naborov.
Poglavje~\ref{sec:diskusija} vsebuje interpretacijo rezultatov, omejitve in možnosti nadaljnjega dela.
Nalogo sklenemo v poglavju~\ref{sec:zakljucek}.

\chapter{Sorodna dela}
\label{sec:sorodna-dela}

To poglavje umešča nalogo med obstoječe pristope za delo s podatki sledilnika pogleda.
Sorodna dela v tem poglavju obravnavamo z vidika vprašanja, kako podatke sledilnika pogleda predstaviti za klasifikacijo čustvenih oznak in kakšno klasifikacijsko uspešnost pri tem dosegajo.
Najprej predstavimo grafovske nevronske mreže za časovne vrste in časovno-prostorske podatke, saj je osrednji del naloge pretvorba časovnega okna meritev sledilnika pogleda v graf.
Na koncu povzamemo še učenje predstavitev iz signala sledilnika pogleda ter položaj te naloge glede na sorodna dela.

\section{Grafovske nevronske mreže za časovno-\linebreak prostorske podatke}

Za modeliranje časovnih vrst se pogosto uporabljajo rekurentne nevronske mreže, kot so mreže LSTM, in transformerji, ki modelirajo časovne odvisnosti med meritvami~\cite{jin2024gnn4ts}.
Grafovska predstavitev je smiselna, kadar želimo poleg časovne urejenosti eksplicitno modelirati tudi odnose med več časovnimi vrstami ali prostorske odnose med meritvami~\cite{jin2024gnn4ts}.
Pregled GNN za časovne vrste to področje obravnava skozi naloge napovedovanja, klasifikacije, imputacije in zaznavanja anomalij.
Za to nalogo je pomembna predvsem formulacija, pri kateri posamezno časovno okno pretvorimo v grafovsko predstavitev in oznako napovemo na ravni celotnega grafa in s tem torej časovnega okna.

Pri podatkih sledilnika pogleda imajo meritve poleg časovnega vrstnega reda tudi prostorsko komponento, saj je vsaka veljavna meritev vezana na položaj pogleda na zaslonu.
Podobno ločevanje časovnih in prostorskih odvisnosti uporabljajo časovno-prostorske grafovske mreže, na primer ST-GCN pri napovedovanju prometa, kjer grafovska konvolucija modelira prostorske povezave med senzorji, časovna konvolucija pa dinamiko signala skozi čas~\cite{yu2018stgcn}.
V tej nalogi je ta pristop uporaben kot metodološka analogija, saj ne napovedujemo prihodnjih meritev, temveč klasificiramo oznako časovnega okna meritev.

\section{Učenje predstavitev iz signalov}
\label{subsec:sorodna-predstavitve}

Poleg ročno oblikovanih značilk obstajajo pristopi, ki se predstavitev naučijo neposredno iz zaporedij meritev.
S tem zmanjšajo odvisnost od ročne izbire značilk, hkrati pa so njihove predstavitve vezane na podatke in naloge, uporabljene pri predtreniranju.
V nalogi takšni modeli služijo kot most med klasičnimi modeli na ročno agregiranih značilkah in grafovskimi modeli, ki se predstavitev učijo na sami strukturi podatkov.

\ModelName{GazeMAE} je primer odprtokodnega modela, zasnovanega posebej za očesne gibe, ki z dvema časovno-konvolucijskima avtokodirnikoma ločeno kodira zaporedji koordinat pogleda in iz njih izračunanih hitrosti.
Pri samonadzorovanem učenju avtokodirnika rekonstruirata delno izpuščene vrednosti signalov~\cite{bautista2021gazemae}.
\ModelName{MOMENT} je splošnejši prednaučen temeljni model za časovne vrste, ki signal razdeli na dele fiksne dolžine in jih kodira z arhitekturo transformerjev.
Pri predtreniranju naključno zakrije nekatere dele signala in se uči njihove rekonstrukcije iz preostalega konteksta, zato je uporaben kot primerjava pri ostalih signalih~\cite{goswami2024moment}.

\section{Položaj te naloge glede na sorodna dela}

Pregled sorodnih del kaže, da se podatki sledilnika pogleda pri prepoznavanju čustvenih stanj najpogosteje uporabljajo v obliki ročno oblikovanih značilk ali zaporedij meritev~\cite{lim2020emotion,tarnowski2020eye,lu2015combining_eye_eeg}.
Grafovski modeli so pogosteje uporabljeni pri drugih vrstah časovno-prostorskih podatkov oziroma pri fizioloških signalih, predvsem pri EEG, kjer vozlišča običajno predstavljajo elektrode ali možganska področja~\cite{song2020dgcnn,zhong2022rgnn,liu2024gnn_eeg_survey}.
Obstajajo tudi pristopi, ki grafovske modele uporabljajo pri večmodalni prepoznavi čustev iz EEG in podatkov očesnih gibov~\cite{wu2021emotion,lu2015combining_eye_eeg}, vendar nismo našli dela, ki bi obravnavalo izključno podatke sledilnika pogleda kot časovno-prostorski graf posameznih meritev.

Ta naloga se zato osredotoča na primerjavo različnih predstavitev podatkov sledilnika pogleda.
Čustvene oznake uporablja kot evalvacijsko nalogo, s katero primerjamo predstavitve istega časovnega okna: agregirane značilke za klasične modele, naučene predstavitve iz prednaučenih kodirnikov~\cite{bautista2021gazemae} in časovno-prostorsko grafovsko predstavitev za GNN.

Grafovski del naloge preverja, ali eksplicitne časovne, prostorske in fiksacijske relacije med meritvami pogleda prispevajo k uporabnejši predstavitvi okna.
V predlagani predstavitvi vozlišča predstavljajo posamezne meritve, povezave pa časovno bližino, prostorsko bližino in pripadnost isti fiksaciji.
Naloga torej povezuje (1) binarno klasifikacijo čustvenih oznak iz podatkov sledilnika pogleda, (2) GNN za časovno-prostorske podatke ter (3) primerjavo ročno oblikovanih, naučenih in grafovskih predstavitev časovnih oken.

\chapter{Podatki in predobdelava}
\label{sec:podatki}

\section{Podatkovna zbirka \MAHNOB{}}
\label{subsec:podatkovna-zbirka}

V delu uporabljamo podatke iz podatkovne zbirke \MAHNOB{}~\cite{soleymani2012mahnob}.
Gre za večmodalno zbirko, namenjeno raziskavam prepoznavanja čustev in implicitnega označevanja multimedijskih vsebin.
Zbirka vsebuje dva glavna eksperimentalna sklopa: prepoznavanje čustvenih stanj iz odzivov na čustvene videoposnetke in implicitno označevanje na podlagi odzivov na pravilne oziroma napačne oznake pod sliko ali posnetkom.
V tej nalogi uporabljamo samo prvi sklop, to je eksperiment vzbujanja čustev z videoposnetki.

Izvirni članek poroča o 30 rekrutiranih udeležencih, od katerih je bilo v analizi uporabljenih 27.
Udeleženci P9, P12 in P15 so bili izključeni zaradi tehničnih težav oziroma nepopolnega zajema podatkov.
V tej nalogi po vzoru izvirnega članka izključimo iste tri udeležence.
V eksperimentu vzbujanja čustev so udeleženci gledali 20 čustveno vzbujajočih videoposnetkov.
Pred vsakim čustvenim posnetkom je bil prikazan še kratek nevtralni posnetek, odziv nanj pa je bil shranjen ločeno.
Po ogledu posameznega posnetka so udeleženci podali samoocene čustvene oznake, vzburjenosti, valence, dominantnosti in predvidljivosti.

Čeprav je \MAHNOB{} večmodalna zbirka, v tej nalogi uporabljamo izključno signale sledilnika pogleda.
Uporabljeni signali so podrobneje opisani v razdelku~\ref{subsec:vhodni-signali}.
Meritve so bile zajete s sledilnikom Tobii X120 pri frekvenci 60~Hz.

\section{Predhodni poskus z zbirko \ESEED{}}
\label{subsec:eseed}

Pred prehodom na \MAHNOB{} je bil del razvoja izveden na podatkovni zbirki \ESEED{}.
Ta del ni osrednji eksperimentalni prispevek naloge, temveč služi kot zgodovinska motivacija za spremembo podatkovne zbirke.
Začetni poskusi so pokazali težave z zanesljivostjo oznak, kakovostjo signalov in stabilnostjo učenja.
Zato smo glavno eksperimentalno evalvacijo izvedli na podatkovni zbirki \MAHNOB{}, ki ima za naš primer uporabnejšo množico signalov sledilnika pogleda.

\section{Ciljne oznake}
\label{subsec:ciljne-oznake}

Oznake v tej nalogi izhajajo iz udeleženčevih samoocen čustvenega razreda po ogledu videoposnetka.
Diskretne čustvene razrede preslikamo v razrede valence po preslikavi iz izvirnega članka~\cite{soleymani2012mahnob}, ki je prikazana v dodatku~\ref{app:preslikava-valence}.
S tem najprej dobimo tri razrede: neprijetno, nevtralno in prijetno valenco.
Nevtralni razred odstranimo, da ciljna naloga vsebuje izrazitejši kontrast med razredoma in s tem modelom omogočimo boljše pogoje za učenje vzorcev v podatkih.

Oznaka valence je določena na ravni kombinacije udeleženca in videoposnetka.
Vsa časovna okna in s tem vsi grafi istega ogleda videoposnetka zato prejmejo isto oznako.

Vzburjenosti v tej nalogi ne klasificiramo.
Razlog je predvsem večja šumnost oznak v uporabljeni podatkovni zbirki.
Kot kaže matrika zmede v dodatku~\ref{app:ujemanje-oznak-s-samoocenami}, se samoocene stopnje vzburjenosti slabše ujemajo z definiranimi razredi kot samoocene valence.


\begin{figure}[t]
    \centering
    \includegraphics[width=0.82\textwidth]{porazdelitev_razredov_stevila_retain_2026-06-12.png}
    \caption{Porazdelitev 10-sekundnih oken po razredih binarne valence.}
    \label{fig:porazdelitev-valenca-stevila}
\end{figure}

\END_NOTE{Slika je delovna verzija iz runa RETAIN_2026-06-12_16-29-08. Pred oddajo preveri, ali je generirana iz končnega protokola, in jo po potrebi regeneriraj v paleti B.}

\section{Priprava vhodnih signalov}
\label{subsec:vhodni-signali}

\subsection{Skupine signalov}

Tukaj navedemo konkretne skupine signalov, ki jih uporabimo v eksperimentih.
V uporabljeni podatkovni zbirki imamo na voljo pet osnovnih vrst signalov. To so:
\begin{itemize}
    \item časovni žig meritve $t_i$,
    \item koordinati pogleda na zaslonu $(x_i,~y_i)$ -- povprečje med meritvama levega in desnega očesa,
    \item velikost leve in desne zenice $(p_i^{(l)},~p_i^{(r)})$,
    \item oddaljenost oči od sledilnika pogleda $d_i$ ter
    \item informacija o fiksacijah, tj. pripadnost fiksaciji ter trajanje fiksacije.
\end{itemize}

Večina podatkovnih zbirk s podatki sledilnika pogleda vsebuje bodisi le koordinate pogleda bodisi le velikosti zenic bodisi oboje.
Nekatere zbirke pa ponujajo tudi dodatne signale, kot sta v našem primeru oddaljenost od sledilnika ter informacije o fiksacijah.
Da bi s primerjalno študijo pokrili čim širši razpon različnih tipov zbirk, \MAHNOB{} podatke razdelimo v štiri delno prekrivajoče se eksperimentalne množice signalov.
V vseh štirih množicah uporabimo tudi informacijo o času.
Množica \emph{samo pogled} poleg časa vsebuje koordinati pogleda, množica \emph{samo zenici} vsebuje velikosti leve in desne zenice, množica \emph{pogled in zenici} združi koordinati pogleda in velikosti zenic, množica \emph{vsi signali} pa doda še oddaljenost od sledilnika pogleda in informacijo o fiksacijah.
Podrobna povezava med temi množicami, značilkami vozlišč in razpoložljivimi tipi povezav je podana pri opisu grafovske konstrukcije v razdelku~\ref{subsec:grafovska-konstrukcija-po-mnozicah}.
\END_NOTE{Z zunanjo pomočjo mentorja ali slovenista najdi/potrdi izraze in jih potem uporabljaj konsistentno v diplomi. Posebej preveri, ali naj bodo končna imena množic \emph{samo pogled}, \emph{samo zenici}, \emph{pogled in zenici} in \emph{vsi signali}, ter nato ta imena uporabi enako v besedilu, tabelah in slikah.}


\subsection{Čiščenje podatkov}
\label{subsec:ciscenje}

Podatke pred uporabo očistimo.
Najprej iz lokalne kopije zbirke \MAHNOB odstranimo podatke, ki jih v tej nalogi ne uporabljamo, in ohranimo le podatke eksperimenta vzbujanja čustev.
Nato odstranimo vrstice brez signalov pogleda ali velikosti zenic.
Na voljo imamo tudi signale veljavnosti, na podlagi katerih odstranimo meritve, za katere je sledilnik poročal, da je meritev za vsaj eno oko neveljavna.
Dodatno uporabimo filter osamelcev \texttt{q01}-\texttt{q99}, s katerim odstranimo izrazite ekstreme, ki izkrivljajo porazdelitve signalov.
Na koncu kratke vrzeli v numeričnih signalih interpoliramo, signala velikosti leve in desne zenice pa dodatno zgladimo s kratkim drsečim povprečjem.

\section{Normalizacija}
\label{subsec:normalizacija}

Posamezne skupine signalov so izražene v različnih velikostnih redih.
Modeli strojnega učenja, kot so \ModelName{SVM}, \ModelName{MLP} ter GNN, so občutljivi na merilo značilk.
Zato signale standardiziramo.
Povprečje in standardni odklon izračunamo na učni podmnožici, nato pa na validacijski in testni podmnožici standardizacijo uporabimo z istimi parametri kot na učni.
Pri grafovskih modelih na enak način standardiziramo tudi značilke povezav, kadar jih uporabljena množica signalov vsebuje.

\section{Segmentacija v časovna okna}
\label{subsec:windowing}

Časovno vrsto razdelimo na neprekrivajoča se časovna okna dolžine $T=10~s$.
Vsako uporabno časovno okno predstavlja en učni primer.
Za $10~s$ smo se odločili na podlagi sorodnega dela na podatkih sledilnika pogleda, kjer so se daljša okna med dolžinami $1$, $3$, $5$ in $10~s$ praviloma izkazala za bolj informativna~\cite{bozak2024feature}.
Daljših oken ne uporabimo, ker bi težje predpostavili, da celotno okno opisuje enoten čustveni odziv, saj se lahko ta med ogledom posnetka že opazno spremeni~\cite{zhang2020rcea}.
Tudi sorodna dela na prepoznavanju čustev iz očesnih signalov pogosto uporabljajo krajša okna, na primer $1$-$2~s$~\cite{wang2018automatic,wu2021emotion}.

Oznaka okna je določena s parom udeleženec-posnetek, zato so vsa okna istega ogleda videoposnetka označena z isto čustveno oznako.
Uporabnost okna določimo glede na kakovost podatkov v oknu.

Pri vzorčni frekvenci 60~Hz bi 10-sekundno okno brez manjkajočih vrednosti vsebovalo 600 vzorcev.
Zaradi čiščenja, odstranjevanja neveljavnih vzorcev in filtriranja je dejansko število vzorcev v oknih manjše.
Povprečno število meritev v oknu po čiščenju znaša \TODO{vstavi številko}.

Po tvorbi časovnih oken se cevovod razdeli glede na uporabljeno predstavitev.
Negrafovski osnovni modeli iz okna dobijo agregirane tabelarične značilke, zamrznjeni prednaučeni predstavitveni modeli neposredno prejmejo časovne vrste izbranih signalov sledilnika pogleda, grafovski modeli pa grafovsko predstavitev okna.
Podrobnosti teh predstavitev so opisane v poglavjih~\ref{sec:grafovska-predstavitev} in~\ref{sec:modeli}.
V tem poglavju opišemo le skupne korake priprave podatkov.

\section{Povzetek predobdelave}

Tabela~\ref{tab:mahnob-predobdelava} povzema, kako se skozi opisane korake predobdelave zmanjša količina podatkov, uporabljena v končni eksperimentalni primerjavi.

\begin{table}[t]
    \centering
    \caption{Zmanjševanje količine podatkov od izvirne zbirke \MAHNOB{} do končnih časovnih oken, uporabljenih v eksperimentalnem cevovodu.}
    \label{tab:mahnob-predobdelava}
    \small
    \begin{tabularx}{\textwidth}{@{}>{\raggedright\arraybackslash}p{0.28\textwidth} c c >{\raggedright\arraybackslash}X@{}}
        \toprule
        Korak & Vrstice & Udeleženci & Opis \\
        \midrule
        Izbira izvirne zbirke \MAHNOB{} & -- & 30 & Večmodalna zbirka; izvirni članek v analizi uporabi 27 udeležencev. \\
        Izbira podmnožice s čustvenimi oznakami & 3\,797\,165 & 24 & Ohranimo le posnetke iz eksperimenta vzbujanja čustev, ki imajo podatke sledilnika pogleda in čustvene oznake. \\
        Izključitev "slabih" udeležencev & 3\,535\,842 & 22 & Izključeni P9, P12 in P15. \\
        Izbira veljavno označenih vrstic & 2\,995\,632 & 22 & Ohranjeni so samo označeni podatki, uporabljeni za učenje modelov čustvenih stanj. \\
        Priprava podatkov za učenje & 2\,814\,005 & 22 & Odstranjene so vrstice brez signalov pogleda ali zenic. \\
        Filter osamelcev \texttt{q01}-\texttt{q99} & 2\,639\,048 & 22 & Robustno filtriranje štirih osnovnih signalov pogleda in zenic. \\
        Tvorba uporabnih 10-sekundnih časovnih oken & 5\,158 & 22 & Vsako okno postane en učni primer oziroma graf. \\
        \bottomrule
    \end{tabularx}
\end{table}

\chapter{Grafovska predstavitev podatkov sledilnika pogleda}
\label{sec:grafovska-predstavitev}

V prejšnjem poglavju smo meritve sledilnika pogleda razdelili na 10-sekundna časovna okna.
V tem poglavju opišemo, kako posamezno okno pretvorimo v graf, ki ga uporabimo kot vhod za grafovske modele.
Vozlišča predstavljajo veljavne meritve znotraj okna, povezave pa izbrane odnose med meritvami.

Najprej podamo najbogatejšo grafovsko predstavitev učnega vzorca, ki uporablja vse razpoložljive signale.
Nato pokažemo, kako se ista osnovna konstrukcija prilagodi štirim množicam signalov iz poglavja~\ref{subsec:vhodni-signali}.
Kadar določena množica signalov ne vsebuje informacije, potrebne za posamezno značilko ali relacijo, ta del predstavitve izpustimo.
Graf zato v vseh primerih predstavlja isto vrsto učnega vzorca, razlikuje pa se po količini vhodne informacije in po razpoložljivih relacijah med meritvami.

\section{Vozlišča}
\subsection{Definicija}

Vsako vozlišče predstavlja eno meritev sledilnika pogleda znotraj časovnega okna:
\begin{equation}
    v_i \leftrightarrow \mathbf{x}_i.
\end{equation}
Če okno vsebuje $n$ veljavnih meritev, potem graf vsebuje $n$ vozlišč.
Vrednost $n$ se lahko med okni razlikuje, saj po čiščenju in filtriranju niso vsa okna popolna.
Grafovska predstavitev naravno podpira spremenljivo število vozlišč.

\subsection{Vektor značilk vozlišča}
Vektor značilk vozlišča za končni graf za množico \emph{vsi signali} je
\begin{equation}
    \mathbf{x}_i =
    \left[
        t_i,
        x_i,
        y_i,
        p_i^{(l)},
        p_i^{(r)},
        d_i,
        f_{dur,i}
    \right],
\end{equation} % fact check the equation
kjer je $t_i$ normaliziran časovni žig za vozlišče $i$, $x_i$ in $y_i$ sta koordinati pogleda, $p_i^{(l)}$ in $p_i^{(r)}$ velikosti leve in desne zenice, $d_i$ je razdalja med sledilnikom in udeleženčevimi očmi, $f_{dur,i}$ pa trajanje fiksacije.

Za ostale podgrafe se vektor značilk ustrezno filtrira.
V vseh primerih ohranimo časovno značilko $t_i$, ostale podatke pa ohranimo glede na obravnavano množico signalov.

\section{Relacije}
Definiramo tri relacije med meritvami: časovno, prostorsko in fiksacijsko.
Posledično v grafu oblikujemo štiri tipe povezav.
Pri tem sta časovna tipa povezav usmerjena -- časovne povezave naprej in nazaj.
Prostorske in fiksacijske povezave pa obravnavamo kot neusmerjene.

\subsection{Zakaj ločimo relacije}

Vsaka povezava ima določen tip relacije, saj časovna povezava, prostorska povezava in fiksacijska povezava opisujejo različne odnose med meritvami.
Ta zapis omogoča, da kasnejše stopnje arhitekturne lestvice relacije obravnavajo ločeno.
Osnovni GCN uporabi isti nabor povezav, vendar relacij modelno ne loči; v tem primeru dobimo enostaven homogen graf.
Podrobnosti o uporabi relacij v posameznih modelih so opisane v poglavju~\ref{sec:modeli}.

\subsection{Časovne povezave}
Časovne povezave povezujejo meritve zaporedne v času.
Naj bodo vozlišča v posameznem oknu urejena po času, tako da je $v_i$ $i$-ta meritev v časovnem zaporedju.

Ločimo usmerjene časovne povezave naprej in nazaj:
\begin{align}
    E_t^+ &= \{(v_i, v_j) : 1 \leq j - i \leq k_t\}, \\
    E_t^- &= \{(v_i, v_j) : 1 \leq i - j \leq k_t\}.
\end{align}

Parameter $k_t \in \mathbb{N}$ določa število časovnih sosedov na vsako stran.
Trenutna konfiguracija uporablja $k_t=1$, zato je vozlišče $v_i$ povezano samo z neposredno predhodno in neposredno naslednjo meritvijo, kadar ti vozlišči obstajata.

Povezave v $E_t^+$ potekajo od prejšnjih proti kasnejšim meritvam, povezave v $E_t^-$ pa od kasnejših proti prejšnjim meritvam.
Pri grafu z $n$ vozlišči to pri $k_t=1$ pomeni $n-1$ časovnih povezav naprej in $n-1$ časovnih povezav nazaj.
Za polno 10-sekundno okno s pričakovanimi $600$ meritvami je to $599$ povezav v vsako smer.


\subsection{Prostorske povezave}
\label{subsec:prostorske-povezave}

Prostorske povezave so neusmerjene povezave, ki povezujejo meritve s podobnim položajem pogleda znotraj istega časovnega okna.
Te povezave lahko definiramo samo za grafe, ki imajo prostorsko informacijo, torej signala $(x,~y)$
Za vsako vozlišče poiščemo $k_s$ najbližjih sosedov v prostoru koordinat pogleda $(x,~y)$.

Prostorsko razdaljo med meritvama $i$ in $j$ definiramo evklidsko:
\begin{equation}
    d_{ij}^{(s)} =
    \sqrt{(x_i - x_j)^2 + (y_i - y_j)^2}.
\end{equation}

Za vsako vozlišče $v_i$ poiščemo $k_s$ najbližjih sosedov v prostoru koordinat pogleda $(x,~y)$ z uporabo k-d drevesa. Če je $v_j$ med temi sosedi, dodamo povezavo med vozliščema $v_i$ in $v_j$.
Ker prostorske povezave obravnavamo kot neusmerjene, povezavo predstavimo z dvema usmerjenima povezavama $(v_i,v_j)$ in $(v_j,v_i)$.
Na koncu odstranimo podvojene povezave znotraj relacije $E_s$.

Število prostorskih povezav zato ni določeno samo z $n$ in $k_s$, ampak tudi z geometrijo položajev pogleda.
Pred simetrizacijo vsako vozlišče prispeva $k_s$ sosedskih izbir, po simetrizaciji in odstranitvi podvojenih parov pa za usmerjeni zapis neusmerjene relacije velja
\begin{equation}
    n k_s \leq |E_s| \leq 2 n k_s.
\end{equation}
Dejansko gostoto prostorskih povezav zato poročamo z opazovanimi statistikami grafov.

\subsection{Fiksacijske povezave}
\label{subsec:fiksacijske-povezave}

Fiksacijske povezave povezujejo meritve, ki pripadajo isti zaznani fiksaciji.
Pripadnost fiksaciji določimo z identifikatorjem fiksacije, podanim v \MAHNOB podatkih.
Ta podatek imamo zgolj v množici \emph{vsi signali}, zato so te povezave vključene le v končni graf.

\subsubsection{Redčenje povezav znotraj fiksacije}

Znotraj posamezne fiksacije ne povežemo vseh vozlišč med seboj, saj bi to povzročilo preveliko število fiksacijskih povezav.
Namesto tega uporabimo redkejšo konstrukcijo, ki ohrani informacijo o pripadnosti isti fiksaciji, hkrati pa zmanjša prostorsko in časovno zahtevnost konstrukcije grafa ter učenja na grafu.
To konstrukcijo sami imenujemo \emph{razširjene povezave znotraj fiksacije} (angl. \emph{dilated intra-fixation edges}).
Fiksacijo obravnavamo kot neprekinjeno zaporedje vozlišč z istim veljavnim identifikatorjem \texttt{fixation-index}.
Če ima taka fiksacija $F$ vozlišč, jim znotraj fiksacije priredimo lokalne indekse $0,\ldots,F-1$.

Gostoto povezav uravnava parameter $k_f$, ki določa število fiksacijskih povezav, ki jih generira vsako vozlišče.
V trenutni konfiguraciji uporabljamo $k_f=3$.
Za posamezno fiksacijo najprej izračunamo korak razširitve
\begin{equation}
    s = \max(1, \lfloor F/k_f + 0{,}5 \rfloor).
\end{equation}
Množico kandidatnih odmikov nato definiramo kot
\begin{equation}
    D = \{(1 + q s) \bmod F \mid q = 0,\ldots,k_f-1\}.
\end{equation}
Morebitni odmik $0$ in podvojene odmike odstranimo.
Za vsak odmik $d \in D$ dodamo neusmerjene povezave
\begin{equation}
    \{v_i, v_{(i+d)\bmod F}\},
    \qquad i = 0,\ldots,F-1.
\end{equation}
Pri zapisu grafa za posredovanje sporočil vsako neusmerjeno povezavo predstavimo z dvema usmerjenima povezavama.
Podvojene povezave znotraj relacije $E_f$ odstranimo.

Na primer, pri fiksaciji dolžine $F=30$ in parametru $k_f=3$ dobimo $s=10$.
Vozlišče z lokalnim indeksom $0$ torej povežemo z vozlišči $1, 11$ in $21$.
Odmik $1$ ohrani časovno lokalno komunikacijo znotraj fiksacije, večji odmiki pa omogočijo širjenje informacij tudi med časovno bolj oddaljenimi meritvami iste fiksacije.

Razširjene povezave imajo večji doseg kot zgolj zaporedne fiksacijske povezave, hkrati pa se izognejo kombinatorični eksploziji, ki bi nastala pri polnem kliku.
Pri polnem kliku bi namreč število povezav naraščalo kvadratno z dolžino fiksacije.
Pri analiziranih 10-sekundnih oknih dobimo približno $1.956$ fiksacijskih povezav na graf pri $k_f=2$, $2.908$ pri $k_f=3$, $4.163$ pri $k_f=5$ in $7.273$ pri $k_f=10$.
Pri privzeti izbiri $k_f=3$ so fiksacijske povezave predstavljale približno $45{,}7\,\%$ vseh povezav v grafu (relacije \texttt{temporal}, \texttt{spatial} in \texttt{fixation}).
%TODO: fact check številke

Izbira $k_f=3$ je zato kompromis med dosegom in gostoto grafa.
Ker vsaka plast GNN omogoča izmenjavo informacij med neposredno povezanimi vozlišči, se doseg znotraj fiksacije povečuje tudi z večanjem števila plasti. % omeni kako se povečuje (če je to enostavno reči) -- kvadratno, eksponentno, ...?
Pri $k_f=3$ vozlišče že v eni plasti prejme informacije iz več delov iste fiksacije, večplastni model pa lahko te informacije dodatno propagira po fiksacijski relaciji.
Večje vrednosti $k_f$ povečajo gostoto grafa, računsko zahtevnost in možnosti za prekomerno glajenje, manjše vrednosti pa graf dodatno razredčijo.

\subsection{Značilke povezav}
\label{subsec:utezi-povezav}

Povezavam poleg tipa relacije priredimo tudi značilke povezav, ki opisujejo odnos med izvorno meritvijo $v_i$ in ciljno meritvijo $v_j$.
Osnovne časovne značilke so normalizirani čas izvorne in ciljne meritve, $t_i$ in $t_j$, ter časovna razlika $\Delta t = t_j - t_i$.
Pri časovnih povezavah dodamo še smer $\rho_{ij}$, kjer ima povezava naprej vrednost $\rho_{ij}=1$, povezava nazaj pa $\rho_{ij}=-1$.
Značilke povezav so del najbogatejše grafovske predstavitve; eksplicitno jih uporabi utežena stopnja arhitekturne lestvice, opisana v poglavju~\ref{sec:modeli}.

Pri grafih, ki vsebujejo koordinate pogleda, povezavam dodamo tudi prostorske značilke.
To sta razliki v koordinatah pogleda, $\Delta x = x_j - x_i$ in $\Delta y = y_j - y_i$, ter evklidska razdalja med položajema pogleda na zaslonu $d_{ij}^{screen}$.
Te značilke so zato na voljo pri vseh množicah signalov, razen pri množici \emph{samo zenici}.

Pri množici \emph{vsi signali} povezavam dodamo še značilko, povezano z oddaljenostjo od sledilnika pogleda.
Značilka $\Delta d^{eye}_{ij}$ opisuje razliko v oddaljenosti med ciljno in izvorno meritvijo.

Pregled vseh značilk povezav in njihove uporabe je podan v tabeli~\ref{tab:znacilke-povezav}.

\begin{table}[H]
    \centering
    \caption{Značilke povezav v grafovski predstavitvi.}
    \label{tab:znacilke-povezav}
    \begin{tabularx}{\textwidth}{@{}>{\raggedright\arraybackslash}p{0.18\textwidth} >{\raggedright\arraybackslash}X >{\raggedright\arraybackslash}p{0.30\textwidth}@{}}
        \toprule
        Značilka & Opis & Uporaba \\
        \midrule
        $t_i$ & Normaliziran čas izvorne meritve znotraj časovnega okna. & Vse povezave. \\
        $t_j$ & Normaliziran čas ciljne meritve znotraj časovnega okna. & Vse povezave. \\
        $\Delta t$ & Časovna razlika med ciljno in izvorno meritvijo, $\Delta t = t_j - t_i$. & Vse povezave. \\
        $\rho_{ij}$ & Smer časovne povezave; pri povezavi naprej je $\rho_{ij}=1$, pri povezavi nazaj pa $\rho_{ij}=-1$. & Samo časovne povezave. \\
        $\Delta x$ & Razlika v vodoravni koordinati pogleda, $\Delta x = x_j - x_i$. & Grafi s koordinatami pogleda. \\
        $\Delta y$ & Razlika v navpični koordinati pogleda, $\Delta y = y_j - y_i$. & Grafi s koordinatami pogleda. \\
        $d_{ij}^{screen}$ & Evklidska razdalja med položajema pogleda na zaslonu. & Grafi s koordinatami pogleda. \\
        $\Delta d^{eye}_{ij}$ & Razlika v oddaljenosti od sledilnika pogleda med ciljno in izvorno meritvijo. & Samo \emph{vsi signali}. \\
        \bottomrule
    \end{tabularx}
\end{table}


\section{Grafovska konstrukcija po množicah signalov}
\label{subsec:grafovska-konstrukcija-po-mnozicah}

Za vsako množico signalov iz tabele~\ref{tab:mnozice-signalov} zgradimo graf po isti osnovni konstrukciji.
Razlika med grafi je v značilkah vozlišč in relacijah, ki jih izbrana množica signalov omogoča.
Množica \emph{vsi signali} je najbogatejša predstavitev. 
Ostale množice uporabijo ustrezne podmnožice značilk in povezav.

\begin{table}[H]
    \centering
    \caption{Grafovska konstrukcija po množicah signalov.}
    \label{tab:grafovska-konstrukcija-po-mnozicah}
    \begin{tabularx}{\textwidth}{@{}>{\raggedright\arraybackslash}p{0.20\textwidth} >{\raggedright\arraybackslash}X >{\raggedright\arraybackslash}X@{}}
        \toprule
        Množica signalov & Značilke vozlišč & Tipi povezav \\
        \midrule
        \emph{samo pogled} & čas in koordinati pogleda & časovne naprej, časovne nazaj, prostorske \\
        \emph{samo zenici} & čas in velikosti obeh zenic & časovne naprej, časovne nazaj \\
        \emph{pogled in zenici} & čas, koordinati pogleda in velikosti obeh zenic & časovne naprej, časovne nazaj, prostorske \\
        \emph{vsi signali} & čas, koordinati pogleda, velikosti obeh zenic, oddaljenost od sledilnika pogleda in fiksacijska informacija & časovne naprej, časovne nazaj, prostorske, fiksacijske \\
        \bottomrule
    \end{tabularx}
\end{table}

Slika~\ref{fig:konstrukcija-grafa-primer} prikazuje primer najbogatejše konstrukcije na majhnem izseku časovnega okna za množico \emph{vsi signali}.
Prikaz združi vozlišča, časovne, prostorske in fiksacijske povezave v isti graf.

\begin{figure}[H]
    \centering
    \includegraphics[width=\textwidth]{gnn_graph_far_02_P10_90_avi_t106_86-107_95_kt1_ks1_kf3_f3_n9/vse_povezave.png}
    \caption{Primer grafovske konstrukcije za množico \emph{vsi signali}. Vozlišča predstavljajo meritve v izseku časovnega okna, povezave pa prikazujejo časovne, prostorske in fiksacijske relacije.}
    \label{fig:konstrukcija-grafa-primer}
\end{figure}

Ločeni prikazi posameznih relacij so zbrani v dodatku~\ref{app:grafi}.

\section{Statistika konstruiranih grafov}
\label{subsec:statistika-grafov}

% Namen sekcije:
% - Zaključiti poglavje z lastnostmi grafov, ki dejansko nastanejo iz podatkov.
% - To ni razdelek o zahtevnosti modelov.
% Vsebina:
% - Preračunaj statistiko za končno konfiguracijo in za vse štiri množice signalov.
% - Tabela naj verjetno vsebuje stolpce:
%   - \emph{samo pogled};
%   - \emph{samo zenici};
%   - \emph{pogled in zenici};
%   - \emph{vsi signali}.
% - Vrstice naj bodo na primer:
%   - mediana števila vozlišč;
%   - povprečno število vozlišč;
%   - mediana časovnih povezav naprej;
%   - mediana časovnih povezav nazaj;
%   - mediana prostorskih povezav;
%   - mediana fiksacijskih povezav;
%   - mediana skupnega števila povezav.
% - Pri relacijah, ki v določeni množici signalov ne obstajajo, uporabi nič; pojasni v opombi pod tabelo.
% - Preveri skladnost s končnimi parametri `kt`=1, `ks`=1, `kf`=3.
% Uporabi iz stare verzije:
% - Staro tabelo kot oblikovni začetek, ne kot vir končnih številk.
% Ne uporabi iz stare verzije:
% - Odseka "Računska zahtevnost".
% - Primerjave števila parametrov, časa učenja, časa inference ali modelne zahtevnosti; to sodi v poglavje 7, 8 ali dodatek s hiperparametri.

\chapter{Primerjani modeli}
\label{sec:modeli}

V nalogi se osredotočamo na primerjavo štirih družin modelov strojnega učenja za klasifikacijo čustvenih oznak iz podatkov sledilnika pogleda.
V tem poglavju razložimo, katere so te družine in na kratko opišemo vsak uporabljen model znotraj družine.
Grafovske modele opišemo podrobneje, saj predstavljajo metodološko jedro naloge, poleg tega pa so za podatke sledilnika pogleda manj standardiziran pristop kot uveljavljeni klasični modeli.

Vsi modeli uporabljajo iste izvorne podatke, ista 10-sekundna okna in iste delitve na učne, validacijske in testne množice.
Razlika je v predstavitvi vhodnih podatkov.
Negrafovski osnovni modeli prejmejo agregirano tabelarično predstavitev istega okna, zamrznjeni prednaučeni predstavitveni modeli dobijo časovne vrste izbranih signalov, grafovski modeli pa delujejo na grafu, ki predstavlja časovno okno meritev.

\section{Pregled primerjanih pristopov}
\label{subsec:pregled-primerjanih-pristopov}

V tabeli~\ref{tab:pregled-primerjanih-modelov} je prikazano, katere modele uporabljamo, kateri družini vsak model pripada, kakšne podatke sprejme na vhod in kakšno vlogo ima v naši primerjalni študiji.

\begin{table}[!htbp]
    \centering
    \caption{Pregled primerjanih modelov in njihovih vhodnih predstavitev.}
    \label{tab:pregled-primerjanih-modelov}
    \footnotesize
    \setlength{\aboverulesep}{0.25ex}
    \setlength{\belowrulesep}{0.4ex}
    \begin{tabularx}{\textwidth}{@{}>{\raggedright\arraybackslash}m{0.22\textwidth} >{\raggedright\arraybackslash}m{0.23\textwidth} >{\raggedright\arraybackslash}m{0.20\textwidth} >{\raggedright\arraybackslash}X@{}}
        \toprule
        \textbf{Družina modelov} & \textbf{Model} & \textbf{Vhod v model} & \textbf{Vloga v primerjavi} \\
        \midrule
        \multirow[c]{6}{=}{\raggedright izhodiščna klasifikatorja} & Naključni & \multirow[c]{6}{=}{\raggedright brez učenja iz značilk} & spodnja meja naključnega ugibanja \\
        \cmidrule(lr){2-2}\cmidrule(l){4-4}
        & Večinski & & spodnja meja napovedovanja najpogostejšega razreda \\
        \cmidrule{1-4}
        \multirow[c]{3.6}{=}{\raggedright klasični modeli strojnega učenja} & \ModelName{SVM} & \multirow[c]{3}{=}{\raggedright agregirane značilke okna} & \multirow[c]{3}{=}{\raggedright klasični model strojnega učenja} \\
        \cmidrule(lr){2-2}
        & \ModelName{LightGBM} & & \\
        \cmidrule(lr){2-2}
        & \ModelName{MLP} & & \\
        \cmidrule{1-4}
        \multirow[c]{6}{=}{\raggedright zamrznjeni prednaučeni predstavitveni modeli} & \ModelName{GazeMAE} & koordinati pogleda $(x, y)$ & prednaučena predstavitev za množico \emph{samo pogled} \\
        \cmidrule(l){2-4}
        & \ModelName{MOMENT} & časovne vrste izbranih signalov & prednaučena predstavitev za množice s signali zenic oziroma več signali \\
        \cmidrule{1-4}
        \multirow[c]{10.5}{=}{\raggedright grafovski modeli} & \ModelName{GCN} & homogeni graf okna & osnovni grafovski model \\
        \cmidrule(l){2-4}
        & \ModelName{HeteroGCN-mean} & \multirow[c]{3.5}{=}{\raggedright heterogeni graf okna} & ločena obravnava relacij s preprosto fuzijo \\
        \cmidrule(lr){2-2}\cmidrule(l){4-4}
        & \ModelName{HeteroGCN-MLP} & & ločena obravnava relacij z učljivo fuzijo \\
        \cmidrule(l){2-4}
        & \ModelName{HeteroGCN-MLP-w} & heterogeni graf okna in značilke povezav & najkompleksnejša stopnja arhitekturne lestvice \\
        \bottomrule
    \end{tabularx}
\end{table}


\section{Negrafovski osnovni modeli}
\label{subsec:negrafovski-osnovni-modeli}

\subsection{Izhodiščna klasifikatorja}
\label{subsec:izhodiscna-klasifikatorja}

Večinski in naključni klasifikator predstavljata spodnjo primerjalno mejo.
Večinski klasifikator vedno napove večinski razred učne množice, naključni klasifikator pa razrede napoveduje z enako verjetnostjo ne glede na distribucijo podatkov.
Pri binarni klasifikaciji to pomeni verjetnost $50~\%$ za vsak razred.
Če model ne preseže rezultatov teh dveh klasifikatorjev, težko trdimo, da se iz vhodnih podatkov nauči uporabnih vzorcev.

\subsection{Klasični modeli strojnega učenja}
\label{subsec:klasicni-modeli-strojnega-ucenja}

Za pravično primerjavo negrafovskim osnovnim modelom podamo agregirane značilke istih signalov, kot jih prejme GNN.
V primeru množice \emph{vsi signali} so to koordinati pogleda $(x,~y)$, velikosti leve in desne zenice, oddaljenost od sledilnika pogleda ter informacija o fiksacijah.
Dodamo tudi normaliziran čas meritve znotraj okna.
Za vsako vhodno značilko izračunamo opisne statistike na ravni okna: povprečje, standardni odklon, minimum, maksimum, razpon, mediano, prvi kvartil, tretji kvartil in medkvartilni razmik.
Pri fiksacijah dodamo še število fiksacij, delež meritev v oknu, ki pripadajo fiksacijam, vsoto trajanj, povprečno trajanje in največje trajanje fiksacije.
V primeru množice \emph{vsi signali} je eno okno v tabelarični predstavitvi opisano z 68 agregiranimi značilkami.
Takšna pretvorba časovnega signala v opisne statistike na ravni okna oziroma odziva je običajen pristop pri klasičnih modelih strojnega učenja za podatke pogleda~\cite{bozak2024feature,lim2020emotion,tarnowski2020eye}.

\subsubsection{\ModelName{SVM}}

\ModelName{SVM} uporabimo kot klasični model z RBF jedrom nad agregiranimi značilkami okna~\cite{cortes1995support}.
RBF jedro modelu omogoča nelinearno ločevanje razredov, zato je \ModelName{SVM} uporaben kot močan negrafovski osnovni model.
Ker je \ModelName{SVM} občutljiv na velikostni red vhodnih značilk, agregirane značilke pred učenjem standardiziramo.

\subsubsection{\ModelName{LightGBM}}

\ModelName{LightGBM} uporabimo kot predstavnika gradientno spodbujenih odločitvenih dreves~\cite{ke2017lightgbm}.
Model je primeren za tabelarične podatke, zato je naravna primerjava za agregirane značilke časovnih oken.
Za razliko od \ModelName{SVM} in \ModelName{MLP} standardizacije vhodnih značilk ne potrebuje, saj odločitvena drevesa delujejo na pragovih posameznih značilk.

\subsubsection{\ModelName{MLP}}

Osnovni \ModelName{MLP} uporabimo kot primer osnovnega nevronskega modela~\cite{goodfellow2016deep}.
Vhod modela je vektor agregiranih značilk enega 10-sekundnega okna, izhod pa logiti za razrede ciljne naloge.
Arhitektura je namerno preprosta: dve skriti polno povezani plasti z aktivacijo ReLU in izpuščanjem, za njima pa izhodna linearna plast.
Model učimo z navzkrižno entropijo.

\section{Zamrznjeni prednaučeni predstavitveni modeli}
\label{subsec:zamrznjeni-prednauceni-modeli}

Poleg modelov nad ročno agregiranimi značilkami uporabimo tudi dva prednaučena kodirnika signalov: \ModelName{GazeMAE} za podatke pogleda in \ModelName{MOMENT} za ostale časovne vrste~\cite{bautista2021gazemae,goswami2024moment}.
Takšni modeli vhodnega okna ne pretvorijo v vnaprej določene statistike, temveč iz neposrednega signala izračunajo vektorsko predstavitev.
V tej nalogi ju uporabimo kot zamrznjena kodirnika, zato njunih parametrov med učenjem ne spreminjamo.
Na izračunanih predstavitvah nato učimo ločeno klasifikacijsko glavo \ModelName{MLP}.
Ta \ModelName{MLP} je arhitekturno drugačen, vendar ločen in nepovezan s samostojnim klasifikatorjem \ModelName{MLP} iz družine klasičnih modelov strojnega učenja.

\subsection{\ModelName{GazeMAE}}
\label{subsec:gazemae-model}

Za podatke pogleda uporabimo odprtokodni prednaučeni kodirnik \ModelName{GazeMAE}~\cite{bautista2021gazemae}.
Model prejme samo koordinati pogleda $(x,~y)$, zato ga samostojno uporabimo le za množico \emph{samo pogled}.
Ker je bil \ModelName{GazeMAE} učen na $2~s$ oknih, vzorčenih s frekvenco $500~Hz$, uporabljena $10~s$ okna najprej pretvorimo v obliko, ki jo model pričakuje.
Koordinati pogleda razdelimo na pet neprekrivajočih se $2~s$ kosov in ju prevzorčimo na ciljno frekvenco $500~Hz$.
Kodirnik temelji na časovno-konvolucijskih mrežah, ki ločeno kodirajo položaj in hitrost pogleda.
Iz predstavitev posameznih kosov nato z združevanjem pridobimo predstavitev celotnega okna s~512 razsežnostmi.

\subsection{\ModelName{MOMENT}}
\label{subsec:moment-model}

\ModelName{MOMENT} uporabimo kot splošni prednaučeni model za časovne vrste~\cite{goswami2024moment}.
V nasprotju z \ModelName{GazeMAE} ni vezan samo na koordinate pogleda, zato ga uporabimo za večkanalne časovne vrste signalov, ki pripadajo izbrani množici signalov v posameznem eksperimentu.
To so lahko velikosti zenic, oddaljenost od sledilnika pogleda in trajanje fiksacije.
Samostojno ga uporabimo le za množico \emph{samo zenici}.
Vhodno okno pred uporabo prevzorčimo na 512 časovnih korakov.
\ModelName{MOMENT} temelji na arhitekturi transformerjev za časovne vrste in v uporabljeni nastavitvi izračuna predstavitev okna s~768 razsežnostmi.

\subsection{Združene predstavitve \ModelName{GazeMAE} in \ModelName{MOMENT}}

Pri množicah \emph{pogled in zenici} ter \emph{vsi signali} uporabimo združeno predstavitev obeh kodirnikov.
Predstavitev \ModelName{GazeMAE} in predstavitev \ModelName{MOMENT} konkateniramo v en skupni vektor, ki je nato vhod v klasifikacijsko glavo \ModelName{MLP}.

V rezultatih zamrznjene prednaučene predstavitvene modele obravnavamo kot eno
družino modelov \ModelName{GazeMAE/MOMENT}.
Pri množici \emph{samo pogled} ta vrstica pomeni uporabo \ModelName{GazeMAE}, pri množici
\emph{samo zenici} uporabo \ModelName{MOMENT}, pri množicah \emph{pogled in zenici} ter
\emph{vsi signali} pa združeno predstavitev obeh kodirnikov. Za to smo se odločili zato, da ohranimo tabelo rezultatov pregledno.


\section{Grafovski modeli}
\label{subsec:grafovski-modeli}

Grafovski modeli uporabljajo grafovsko predstavitev okna iz poglavja~\ref{sec:grafovska-predstavitev}.
Pri isti množici signalov vse stopnje uporabljajo iste vozliščne značilke, iste relacije in ista časovna okna.
Razlikujejo se po tem, kako obravnavajo tipe povezav, kako združujejo relacijske predstavitve in ali uporabljajo značilke povezav za izračun naučenih uteži povezav.
Najprej definiramo osnovno GNN, nato pa s tremi nadgradnjami postopno gradimo končni model.
V tabeli~\ref{tab:arhitekturna-lestvica-gnn} so na kratko predstavljene razlike med vsemi štirimi stopnjami grafovskih modelov.

\begin{table}
    \centering
    \caption{Arhitekturna lestvica grafovskih modelov.}
    \label{tab:arhitekturna-lestvica-gnn}
    \small
    \begin{tabularx}{\textwidth}{@{}>{\raggedright\arraybackslash}p{0.24\textwidth} >{\centering\arraybackslash}p{0.10\textwidth} >{\centering\arraybackslash}p{0.12\textwidth} >{\centering\arraybackslash}p{0.10\textwidth} >{\raggedright\arraybackslash}X@{}}
        \toprule
        \textbf{Model} & \textbf{Ločene relacije} & \textbf{Fuzija relacij} & \textbf{Uteži povezav} & \textbf{Namen stopnje} \\
        \midrule
        \ModelName{GCN} & ne & -- & ne & preveri osnovno grafovsko predstavitev v homogenem grafu \\
        \ModelName{HeteroGCN-mean} & da & povprečje & ne & preveri vpliv ločene obravnave relacij \\
        \ModelName{HeteroGCN-MLP} & da & \ModelName{MLP} & ne & preveri vpliv naučene fuzije relacijskih predstavitev \\
        \ModelName{HeteroGCN-MLP-w} & da & \ModelName{MLP} & da & preveri vpliv naučenih predznačenih uteži povezav \\
        \bottomrule
    \end{tabularx}
\end{table}


\subsection{Skupni cevovod grafovskih modelov}
\label{subsec:skupni-deli-grafovske-arhitekture}

Grafovski modeli na vhod prejmejo graf glede na izbrano množico signalov, kot definirano v~\ref{sec:grafovska-predstavitev}.
Razsežnost vhodnih preslikav prilagodimo izbrani množici signalov -- ločeno za značilke vozlišč, pri modelu \ModelName{HeteroGCN-MLP-w} pa tudi za značilke povezav.

Vhodne vektorje najprej preslikamo v skupni latentni prostor z dvoslojnim \ModelName{MLP}.
Tako dobimo začetno predstavitev vsakega vozlišča, ki jo nato posodabljajo grafovske plasti.
Te plasti posredujejo sporočila po povezavah grafa in pri vsakem vozlišču združijo informacije iz njegove okolice.
V grafovskih blokih uporabimo aktivacijo GELU, rezidualne povezave, normalizacijo plasti in izpuščanje.
Ti gradniki niso predmet primerjave med grafovskimi modeli.
Primerjava se osredotoča na obravnavo relacij, fuzijo relacijskih predstavitev in uporabo značilk povezav.
Iz končnih predstavitev vozlišč nato s pozornostnim združevanjem\footnote{
Ker napovedujemo oznako celotnega okna, predstavitve vozlišč združimo v predstavitev grafa.
Uporabimo $\mathbf{h}_G = \sum_{v \in V} \alpha_v \mathbf{h}_v$, kjer je $\mathbf{h}_G$ predstavitev grafa, $\alpha_v$ utež pozornosti za vozlišče $v$, $\mathbf{h}_v$ pa predstavitev vozlišča $v$.
Uteži pozornosti izračunamo s softmax normalizacijo: $\alpha_v = \exp(g(\mathbf{h}_v)) / \sum_{u \in V} \exp(g(\mathbf{h}_u))$, kjer je $g$ naučena funkcija, ki predstavitvi vozlišča priredi skalarno oceno pomembnosti.
} izračunamo predstavitev celotnega grafa, to pa pošljemo skozi klasifikacijsko glavo \ModelName{MLP}.
Glava izračuna dva logita, enega za neprijetno in enega za prijetno valenco.
Napoved je razred z večjim logitom.
Pri poročanju verjetnosti logite normaliziramo s softmaxom.
Podrobnosti o številu plasti, velikosti skritih predstavitev, verjetnosti izpuščanja in ostalih hiperparametrih so opisane v dodatku~\ref{app:hiperparametri}.

Vse GNN treniramo od začetka do konca (ang. \emph{end-to-end}), kar pomeni, da se vsi učljivi deli modela optimizirajo hkrati glede na klasifikacijsko izgubo.
To vključuje začetno preslikavo značilk vozlišč, grafovske plasti, morebitno fuzijo relacij, pozornostno združevanje in klasifikacijsko glavo.
V nadaljevanju opišemo specifike vsake arhitekture v štiri-stopenjski arhitekturni lestvici, začenši z najenostavnejšim -- s homogenim \ModelName{GCN}.

\subsection{\ModelName{GCN}}
\label{subsec:model-gcn}

\ModelName{GCN} je v osnovi homogeni grafovski model.
Uporablja isto grafovsko predstavitev okna, kot je bila opisana v poglavju~\ref{sec:grafovska-predstavitev}, vendar ne ločuje časovnih, prostorskih in fiksacijskih relacij.
Vse povezave združi v en homogen graf, posredovanje sporočil pa izvede z običajnimi plastmi \texttt{GCNConv}.
Model ne uporablja značilk povezav in posledično naučenih uteži povezav.

Tehnično lahko model zapišemo kot:
\begin{equation}
    X
    \rightarrow
        MLP_{\text{vhod}}
    \rightarrow
    f_{\GCN}
    \rightarrow
    \text{READOUT}
    \rightarrow
    f_{\text{head}},
\end{equation}
kjer je $X$ matrika značilk vozlišč, $MLP_{\text{vhod}}$ vhodni dvoslojni perceptron, $f_{\GCN}$ zaporedje grafovskih konvolucijskih plasti, korak $READOUT$ združi predstavitve vozlišč v predstavitev celotnega grafa, $f_{\text{head}}$ pa iz te predstavitve izračuna napoved razreda.
TODO{Zamenjaj to enačbo z diagramom, ko ga imaš. Z diagramom lažje prikažeš že samo razliko med tem in naslednjim modelom in imaš potem majhen diagram v vsaki subsekciji in obkrožiš razlike pri vsaki nadgradnji}

\subsection{\ModelName{HeteroGCN-mean}}
\label{subsec:model-heterogcn-mean}

\ModelName{HeteroGCN-mean} dobimo tako, da \ModelName{GCN} nadgradimo s heterogeno grafovsko predstavitvijo.
Kot opisano v~\ref{sec:grafovska-predstavitev} tukaj ločimo različne tipe relacij.
Relacijske predstavitve vozlišč združimo v skupno predstavitev vozlišča s preprostim povprečenjem.
Namen te stopnje je izolirati vpliv ločevanja relacij.
Kot prikazano v diagramu~\ref{fig:diagram-heterogcn-mean} je torej sprememba glede na homogeni \ModelName{GCN} je zato samo v tem, da se sporočila najprej izračunajo ločeno po relacijah, nato pa združijo s povprečenjem.
\TODO{ko bom naredil diagram, ga dodam tu}

\subsection{\ModelName{HeteroGCN-MLP}}
\label{subsec:model-heterogcn-mlp}

\ModelName{HeteroGCN-MLP} dobimo tako, da \ModelName{HeteroGCN-mean} nadgradimo z naučeno fuzijo relacijskih predstavitev.
Različne predstavitve vozlišč -- tiste, ki so v danem modelu na voljo izmed časovne naprej, časovne nazaj, prostorske ter fiksacijske -- najprej konkateniramo, nato pa jih z dvoslojnim perceptronom združimo v skupno predstavitev vozlišča.
S tem dopuščamo, da model sam določi, kateri deli katere relacije so pomembnejši za končno predstavitev.
Kot prikazano v diagramu~\ref{fig:diagram-heterogcn-mlp} je torej sprememba  glede na \ModelName{HeteroGCN-mean} je torej v zamenjavi ročno določene fuzije z naučeno fuzijo relacijskih predstavitev.
\TODO{ko bom naredil diagram, ga dodam tu}

\subsection{\ModelName{HeteroGCN-MLP-w}}
\label{subsec:model-heterogcn-mlp-w}

\ModelName{HeteroGCN-MLP-w} dobimo tako, da \ModelName{HeteroGCN-MLP} nadgradimo z naučenimi utežmi povezav.
Za izračun uteži povezav ima model tri večslojne perceptrone: $f_w^{\text{čas}}$, $f_w^{\text{prostor}}$ in $f_w^{\text{fiksacija}}$.
Ti na podlagi vektorja značilk posamezne povezave, katerega komponente so povzete v tabeli~\ref{tab:znacilke-povezav}, izračunajo skalarno predznačeno utež, ki določa, kako močno posamezna povezava vpliva na posredovanje sporočil.
Pri posamezni množici signalov se uporabijo le funkcije za relacije, ki so v grafu prisotne.
Časovne povezave naprej in nazaj si delijo isto funkcijo $f_w^{\text{čas}}$, smer povezave pa je podana kot dodatna značilka povezave $\rho_{ij}$.
Formule za izračun in normalizacijo uteži so podane v dodatku~\ref{app:naucene-utezi-povezav}.
Namen te \emph{končne} stopnje je preveriti, ali modelu koristi razlikovati moč in predznak prispevka posameznih povezav znotraj relacije.
\TODO{ko bom naredil diagram, ga dodam tu}

\subsection{Izbira grafovske konvolucije}
\label{subsec:izbira-grafovske-konvolucije}

V vseh grafovskih modelih uporabimo plast \texttt{GCNConv}~\cite{kipf2017gcn}.
\texttt{GCNConv} je enostaven in uveljavljen operator posredovanja sporočil, ki ga lahko ohranimo nespremenjenega skozi celotno arhitekturno lestvico.
Alternativni operatorji, kot sta GraphSAGE ali GIN, bi v našem primeru odprli dodatne arhitekturne odločitve in s tem otežili neposredno primerjavo stopenj.

Izbiro smo preverili tudi z manjšim dodatnim poskusom na množici signalov \emph{pogled in zenici}.
Razlike med preizkušenimi operatorji so bile razmeroma majhne, \texttt{GCNConv} pa je dosegel primerljive in konsistentne rezultate tako pri osnovnem homogenem modelu kot tudi uteženi heterogeni arhitekturi.
Podrobnosti primerjave so podane v dodatku~\ref{app:izbira-grafovske-konvolucije}.

\section{Povzetek poglavja}
\label{subsec:povzetek-modelov}

V delu primerjamo štiri družine modelov različnih kompleksnosti in različnih arhitektur: izhodiščna klasifikatorja, klasične modele strojnega učenja, zamrznjena prednaučena kodirnika in grafovske modele.
Vhodne predstavitve za te modele posledično razdelimo na tri vrste: agregirane tabelarične značilke, neagregirane časovne vrste signalov in grafovske predstavitve podatkov sledilnika pogleda.
Grafovska lestvica štirih arhitektur preverja, kaj prispevajo sama grafovska struktura, ločena obravnava relacij, naučena fuzija relacijskih predstavitev in uteženost povezav.
V naslednjem poglavju opišemo evalvacijski protokol za primerjavo opisanih modelov na štirih množicah signalov.

\chapter{Eksperimentalni protokol}
\label{sec:eksperimenti}

V prejšnjih poglavjih smo definirali podatke, ciljne oznake, časovna okna,
grafovsko predstavitev in primerjane modele.
V tem poglavju določimo, kako te komponente ovrednotimo v skupni eksperimentalni
primerjavi.
Glavna evalvacijska naloga je binarna klasifikacija nizke proti visoki valenci.
Primerjava je zasnovana kot kombinacija modelov iz poglavja~\ref{sec:modeli} in
množic signalov iz razdelka~\ref{subsec:vhodni-signali}.
Vsi modeli uporabljajo iste 10-sekundne učne vzorce in iste delitve na učno,
validacijsko in testno množico.

\section{Zasnova glavne primerjave}
\label{subsec:zasnova-glavne-primerjave}

Cilj glavne primerjave je odgovoriti na raziskovalno vprašanje~\ref{rv:glavno}: \rvGlavnoBesedilo
Zato v primerjavi spreminjamo le način predstavitve okna in model, ki to predstavitev uporabi.

Ciljna oznaka je binarni razred valence, pri kateri model ločuje med nizko in visoko valenco.
Vse štiri množice iz tabele~\ref{tab:grafovska-konstrukcija-po-mnozicah} obravnavamo kot enakovreden del primerjave, saj nas zanima tudi, kako se klasifikacijska uspešnost spreminja z razpoložljivostjo različnih signalov sledilnika pogleda.

V primerjavi nastopajo vsi modeli iz tabele~\ref{tab:pregled-primerjanih-modelov}: izhodiščna klasifikatorja, klasični modeli nad agregiranimi značilkami, zamrznjeni prednaučeni predstavitveni modeli in grafovski modeli.
S tem primerjamo tri različne načine opisa istega časovnega okna: tabelarični povzetek signalov, vektorsko predstavitev časovne vrste in grafovsko predstavitev meritev.

\section{Validacijski protokol}
\label{subsec:validacijski-protokol}

\subsection{7-kratno prečno preverjanje po subjektih}
\label{subsec:subject-7fold}

Glavno primerjavo izvedemo s 7-kratnim prečnim preverjanjem po subjektih.
Pri vsaki ponovitvi so testni subjekti popolnoma ločeni od subjektov, ki so uporabljeni za učenje in validacijo.
Iz učnega dela izberemo eno skupino subjektov za validacijsko množico, preostali subjekti pa sestavljajo učno množico.
Delitve so pri vseh modelih in vseh množicah signalov enake, določene s fiksnim naključnim semenom.

Takšna delitev ocenjuje posploševanje na nove udeležence.
To je za uporabljeno nalogo pomembno, saj so signali sledilnika pogleda močno odvisni od specifik posameznika.
Ker 22 subjektov ni deljivo s sedem, imajo posamezne testne razdelitve različno velikost -- nekatere vsebujejo tri, druge pa štiri testne subjekte.

Validacijska množica se uporablja med učenjem, testna množica pa samo za končno oceno v posamezni ponovitvi.
Pri nevronskih modelih zato najboljšo različico modela izberemo glede na validacijsko izgubo, nato pa njeno delovanje izmerimo na pripadajoči testni množici.

\subsection{Predhodna preverjanja z LOO po subjektih}
\label{subsec:predhodna-loo-preverjanja}

V zgodnejših fazah smo za množico \emph{vsi signali} izvedli tudi LOO po subjektih.
Ker končna primerjava vključuje več modelov in štiri množice signalov, bi polni LOO protokol bistveno povečal število potrebnih eksperimentov.
Zato ga uporabimo kot podporno preverjanje, glavno primerjavo pa izvedemo s 7-kratnim protokolom.
Primerjava obeh validacijskih protokolov za valenco in množico \emph{vsi signali} je podana v dodatku~\ref{app:razsirjeni-rezultati-binarna-valenca}, kjer vidimo, da je 7-kratno prečno preverjanje primerljivo z metodo LOO.

\section{Nastavitve učenja in hiperparametri}
\label{subsec:nastavitve-ucenja}

V tem razdelku navedemo nastavitve učenja, ki vplivajo na izvedbo eksperimentov.
Nevronske modele učimo z zgodnjim ustavljanjem na validacijski izgubi.
Grafovski modeli uporabljajo operator \texttt{GCNConv}, dve grafovski plasti, skrite predstavitve velikosti 64 ter povezave s parametri $k_t=1$, $k_s=1$ in $k_f=3$, kot so definirani v poglavju~\ref{sec:grafovska-predstavitev}.
Podrobni hiperparametri vseh modelov so podani v dodatku~\ref{app:hiperparametri}.

Nastavitve grafovskih modelov izberemo na podlagi predhodnega kompaktnega mrežnega iskanja.
To iskanje je bilo izvedeno na 3-razredni valenci in ga v ~\ref{app:hiperparametri} poročamo samo kot pomožni poskus za izbiro stabilnih nastavitev.

\section{Metrike vrednotenja}
\label{subsec:metrike}

Glavni metriki vrednotenja sta točnost in makro F1.
Za vsak model in vsako množico signalov poročamo povprečje po sedmih testnih razdelitvah.
Točnost uporabimo kot osnovno metriko, ker ima glavna naloga dva enakovredna in globalno skoraj uravnotežena razreda, kot kaže slika~\ref{fig:porazdelitev-valenca-stevila}.

Makro F1 uporabimo kot dopolnilno glavno metriko.
Čeprav sta razreda globalno skoraj uravnotežena, so lahko posamezne testne razdelitve precej neuravnotežene (tudi do razmerja $1:2$).
Makro F1 izračuna F1 ločeno za vsak razred in nato razreda povpreči z enako utežjo.
S tem dopolni točnost predvsem pri modelih, ki lahko dosežejo sprejemljivo točnost, vendar večino primerov napovedujejo v en sam razred.

Poleg glavnih metrik spremljamo tudi uravnoteženo točnost, preciznost, priklic oziroma občutljivost, uteženi F1, AUC, izgubo in matrike zmede.
Te metrike uporabimo predvsem za preverjanje stabilnosti in za podrobnejšo analizo napak.
V glavnem besedilu se osredotočimo na točnost in makro F1, razširjene metrike pa poročamo v dodatku~\ref{app:razsirjeni-rezultati-binarna-valenca}. % če si podoben stavek napisal v poglavju 8, potem ga tu odstrani.

\section{Spremljanje učenja}
\label{subsec:spremljanje-ucenja}

Poleg končnih metrik shranjujemo tudi podatke o poteku učenja.
Pri nevronskih modelih spremljamo učno in validacijsko izgubo, najboljšo epoho, informacijo o zgodnjem ustavljanju, čas epohe in norme gradientov.
Krivulje izgub služijo predvsem preverjanju stabilnosti optimizacije in zaznavanju potencialnega prekomernega prileganja.

Pri grafovskih modelih spremljamo še nekaj diagnostičnih meritev naučenih predstavitev in napovedi.
Med njimi so varianca predstavitev grafa, povprečna kosinusna podobnost predstavitev, razpon logitov in entropija napovednih verjetnosti.
Te meritve uporabimo za preverjanje, ali se med učenjem pojavijo skoraj konstantne napovedi, zaznavanje kolapsa predstavitev ali znakov prekomernega glajenja.
Te diagnostike ne poročamo med rezultati v poglavju~\ref{sec:rezultati}, temveč le v dodatku~\ref{app:diagnostika-prekomerno-glajenje}.

\chapter{Diskusija}
\label{sec:diskusija}

\TODO{
\begin{itemize}
    \item Diskusijo preusmeriti iz vprašanja, ali GNN absolutno zmaga, v vprašanje, kako se modeli obnašajo pri različnih naborih signalov.
    \item Posebej interpretirati razlike med \emph{samo pogled}, \emph{pogled in zenici} in \emph{vsi signali}: ali zenični signal pomaga nad koordinatami pogleda, ali dodatni signali v \emph{vsi signali} pomagajo ali dodajo šum, ter ali GNN bogatejšo strukturo izkoristi bolje kot negrafovski osnovni modeli.
    \item \ModelName{GazeMAE} obravnavati kot močno referenco za množico \emph{samo pogled}; če premaga modele z več signali, to interpretirati kot pomembno ugotovitev o naučenih predstavitvah pogleda.
    \item \emph{Samo zenici} in \ModelName{MOMENT} opisati samo kot prihodnje delo ali podporni poskus, če ne bosta dokončana.
\end{itemize}
}

\section{Ali grafovska predstavitev pomaga?}
\label{subsec:diskusija-graf-pomaga}

\TODO{
\begin{itemize}
    \item ali finalni GNN premaga negrafovske osnovne modele;
    \item če ne premaga: ali kaže smiselno vedenje;
    \item pri kateri nalogi je graf najbolj koristen;
    \item LOO po subjektih v primerjavi z LOO po posnetkih;
    \item primerjava proti \ModelName{GazeMAE};
    \item povezava z ablations;
    \item konservativen zaključek.
\end{itemize}
}

\section{Vloga časovnih in prostorskih povezav}
\label{subsec:diskusija-edge-types}

\TODO{
\begin{itemize}
    \item časovne povezave;
    \item prostorske povezave;
    \item fiksacijske povezave;
    \item ločevanje časovnih povezav naprej in nazaj;
    \item naučene predznačene uteži povezav;
    \item \texttt{distance-avg};
    \item \texttt{fixation-duration};
    \item značilka povezave \texttt{delta\_distance};
    \item kaj to pomeni za signale sledilnika pogleda;
    \item brez kavzalnih trditev.
\end{itemize}
}

\section{Primerjava z \ModelName{GazeMAE}}
\label{subsec:diskusija-gazemae}

\TODO{
\begin{itemize}
    \item ali \ModelName{GazeMAE} premaga negrafovske osnovne modele;
    \item ali \ModelName{GazeMAE} premaga GNN;
    \item samonadzorovano učenje predstavitev kot uporabna smer;
    \item zamrznjena enkoderja position/velocity;
    \item 512-dim okenska predstavitev;
    \item omejitev: predstavitveni model s prenosom;
    \item omejitev: ne polna reprodukcija predtreniranja.
\end{itemize}
}

\section{Stabilnost učenja in notranje predstavitve}
\label{subsec:diskusija-stabilnost}

\TODO{
\begin{itemize}
    \item Povzemi, ali krivulje izgub kažejo stabilno učenje ali prekomerno prileganje.
    \item Interpretiraj varianco in kosinusno podobnost predstavitev v povezavi s prekomernim glajenjem.
    \item Interpretiraj razpon logitov v povezavi z zanesljivostjo oziroma zasičenostjo napovedne glave.
    \item V eni povedi poročaj, da diagnostika iz dodatka~\ref{app:diagnostika-prekomerno-glajenje} kaže informativno sliko učenja: GNN se uči razločevalnih grafovskih predstavitev in pri glavnih signalnih naborih ne kaže jasnega kolapsa oziroma prekomernega glajenja.
    \item Če diagnostika ne pokaže jasnega vzorca, to napiši previdno in brez pretiranih sklepov.
    \item Poveži ugotovitve z arhitekturnimi odločitvami iz poglavja~\ref{subsec:stabilizacija-ucenja}.
\end{itemize}
}

\section{Primerjava z izvirnim \MAHNOB{} pristopom}
\label{subsec:diskusija-mahnob}

\TODO{
\begin{itemize}
    \item izvirni članek: ročno oblikovane značilke pogleda;
    \item izvirni članek: drugačen cevovod izbire značilk;
    \item izvirni članek: \ModelName{SVM};
    \item ta naloga: grafi iz časovnih oken surovejših signalov;
    \item primerjava metodološko motivacijska;
    \item ne neposredna reprodukcija;
    \item posebej omeniti uporabo \texttt{distance-avg} in fiksacijskih informacij kot približanje širšemu naboru signalov iz literature.
\end{itemize}
}

\section{Omejitve}
\label{subsec:omejitve}

Glavne omejitve naloge so:
\begin{itemize}
    \item uporaba ene glavne podatkovne zbirke;
    \item subjektivnost čustvenih oznak;
    \item omejitev na signale sledilnika pogleda brez drugih fizioloških modalnosti;
    \item občutljivost grafovske konstrukcije na hiperparametre;
    \item računska zahtevnost validacije z LOO protokoli;
    \item možnost prekomernega glajenja v grafovskih modelih;
    \item zaradi subject kfold model med treniranjem vidi patterne recordingov, tako da se lahko uči le teh. narediti bi bilo potrebno tudi CV po recordingih.
\end{itemize}

\section{Možnosti nadaljnjega dela}
\label{subsec:future-work}

Možne smeri nadaljnjega dela vključujejo:
\begin{itemize}
    \item učenje povezav namesto ročne konstrukcije grafa;
    \item uporabo grafovskih transformerjev;
    \item samonadzorovano predtreniranje na neoznačenih signalih sledilnika pogleda;
    \item razvoj grafovskega temeljnega modela za podatke pogleda;
    \item razširitev na večmodalne fiziološke podatke;
    \item evalvacijo na več podatkovnih zbirkah.
\end{itemize}

\TODO{
Pri pisanju tega odseka dodaj tudi dve konkretni smeri nadaljnjega dela:
\begin{itemize}
    \item razmislek in implementacijo uporabe mežikov, saj literatura mežike obravnava kot informativen signal pri prepoznavi čustvenih stanj iz podatkov pogleda~\cite{lim2020emotion,tarnowski2020eye};
    \item razširitev grafovske predstavitve s fiksacijskimi meta-vozlišči, torej idejo 3D lift, kjer bi posamezne meritve ostale osnovna vozlišča, fiksacije pa bi bile dodatna vozlišča višje ravni, povezana z meritvami, ki jim pripadajo.
\end{itemize}
}


\appendix
\chapter{Dodatni rezultati}
\label{app:dodatni-rezultati}

\TODO{
\begin{itemize}
    \item Sem dodaj dodatne tabele po foldih, rezultate manj pomembnih modelov in rezultate, ki niso šli v glavno besedilo.
    \item Dodaj rezultate za \textit{subject LOO} in \textit{recording LOO} pri vseh signalnih naborih: pogled, zenici, pogled+zenici in vsi signali.
    \item Dodaj vse izračunane metrike rezultatov, ne samo \textit{accuracy} in makro-F1, ki ostaneta v glavnem besedilu.
    \item Razširjene metrike: balanced accuracy, weighted-F1, AUC, po potrebi precision/recall po razredih in rezultati po posameznih gubah.
    \item Dodatne matrike zmede: recording LOO, binarne naloge, najboljši negrafovski model in primerjave, ki niso nujne za glavno zgodbo.
    \item Dodaj rezultate primerjave različnih tipov grafovskih konvolucij oziroma GNN slojev, če ostanejo zunaj glavnega besedila.
    \item Dodaj rezultate 3-razredne klasifikacije in kratek komentar primerjave z izvirnim člankom podatkovne zbirke MAHNOB-HCI oziroma HCI-tagging.
    \item Dodatne porazdelitve podatkov: porazdelitev razredov po subjektih, po posnetkih, po izvorni čustveni oznaki in za binarni low/high nalogi.
    \item Dodatne slike ujemanja oznak: podrobnejše heatmape Table 6 razredov proti samoocenam in povzetki label-noise analize, če v glavnem besedilu ostane samo ena skupna slika.
    \item UMAP procesiranih podatkov dodaj sem samo, če je slika zanimiva, vendar premalo ključna ali premalo čista za glavno besedilo. Če UMAP ne doda interpretativne vrednosti, ga ne vključuj.
\end{itemize}
}

\chapter{Hiperparametri}
\label{app:hiperparametri}

\section{Hiperparametri modelov}
\label{app:hiperparametri-modelov}

Tabela~\ref{tab:hiperparametri-modelov} povzema glavne hiperparametre modelov v končnih eksperimentih.
Klasični primerjalni modeli uporabljajo agregirane tabelarične značilke, nevronske primerjalne glave uporabljajo bodisi agregirane značilke bodisi zamrznjene predstavitve, grafovski modeli pa uporabljajo grafovsko predstavitev časovnih oken.

\begin{table}
    \centering
    \caption{Glavni hiperparametri modelov v končnih eksperimentih.}
    \label{tab:hiperparametri-modelov}
    \footnotesize
    \setlength{\tabcolsep}{3pt}
    \renewcommand{\arraystretch}{0.9}
    \begin{tabularx}{\textwidth}{@{}>{\raggedright\arraybackslash}p{0.21\textwidth} >{\raggedright\arraybackslash}p{0.20\textwidth} >{\raggedright\arraybackslash}X@{}}
        \toprule
        \textbf{Model oziroma družina} & \textbf{Podatkovna predstavitev} & \textbf{Glavni hiperparametri} \\
        \midrule
        Naključni model & agregirane značilke & enakomerna naključna napoved; seme 42 \\
        Večinski model & agregirane značilke & napoved večinskega razreda na učni množici \\
        \ModelName{SVM} & agregirane značilke & RBF; $C=1{,}0$; \texttt{probability=true}; standardizacija značilk \\
        \ModelName{LightGBM} & agregirane značilke & \texttt{n\_estimators=100}; \texttt{learning\_rate=0.05}; \texttt{max\_depth=5}; seme 42 \\
        \ModelName{MLP} & agregirane značilke & skrita velikost 64; največ 200 epoh; učna stopnja $10^{-3}$; \texttt{dropout=0.1}; mini-serija 128; potrpežljivost 15 \\
        \ModelName{GazeMAE}/\ModelName{MOMENT} glave & zamrznjene predstavitve & skrita velikost 128; največ 100 epoh; učna stopnja $10^{-3}$; \texttt{dropout=0.2}; mini-serija 128; potrpežljivost 15 \\
        Grafovski modeli & grafi časovnih oken & \texttt{GCNConv}; 2 plasti; \texttt{hidden\_channels=64}; pozornostno združevanje; \texttt{dropout=0.1}; učna stopnja $5 \cdot 10^{-5}$; mini-serija 256; največ 350 epoh; potrpežljivost 20 \\
        \ModelName{GCN} & grafi časovnih oken & homogen graf; brez ločenih relacij in značilk povezav \\
        \ModelName{HeteroGCN-mean} & grafi časovnih oken & ločene relacije; združevanje s povprečenjem \\
        \ModelName{HeteroGCN-MLP} & grafi časovnih oken & ločene relacije; združevanje z \ModelName{MLP} \\
        \ModelName{HeteroGCN-MLP-w} & grafi časovnih oken & kot \ModelName{HeteroGCN-MLP}, z naučenimi utežmi povezav \\
        Grafovska konstrukcija & grafi časovnih oken & $k_t=1$; $k_s=1$; $k_f=3$; standardizirane značilke povezav \\
        \bottomrule
    \end{tabularx}
\end{table}


\section{Predhodno mrežno iskanje}
\label{app:mrezno-iskanje}

Nastavitve grafovske konstrukcije smo izbrali po majhnem predhodnem mrežnem iskanju na 3-razredni valenci, glede na izvorne razrede definirane v tabeli~\ref{tab:label-mapping}.
Iskanje je spreminjalo število grafovskih plasti, velikost skritih predstavitev ter parametre konstrukcije $k_t$, $k_s$ in $k_f$.
Primerjali smo več nastavitev števila plasti in skrite razsežnosti, vendar med njimi ni bilo velikih razlik.
Ker je bil najmanjši model med najboljšimi in najstabilnejšimi, hkrati pa je računsko najvarčnejši, smo za končne poskuse izbrali $L=2$ (število skritih plasti GNN) in $d=64$ (razsežnost skritih plasti).
Pri tej arhitekturni nastavitvi smo v mrežnem iskanju spreminjali $k_t \in \{1,2,3\}$, $k_s \in \{1,2,3\}$ in $k_f \in \{1,2,3,5\}$.
Najvišji rezultat v tem pomožnem iskanju je dosegla nastavitev $k_t=3$, $k_s=1$ in $k_f=5$, vendar so bile razlike majhne.
V končnih poskusih smo uporabili redkejšo nastavitev $k_t=1$, $k_s=1$ in $k_f=3$, ker zmanjša gostoto grafa in računsko zahtevnost ter ohrani stabilno delovanje.

% Vir podatkov: GFM-for-eyetracker/results/gnn_v2_valence_grid_search_2x64_phase2/2026-05-28_16-54-08/grid_summary.csv.
\begin{table}[H]
    \centering
    \caption{Povzetek predhodnega mrežnega iskanja hiperparametrov grafovske konstrukcije.}
    \label{tab:mrezno-iskanje-gnn}
    \footnotesize
    \begin{tabularx}{\textwidth}{@{}>{\raggedright\arraybackslash}p{0.26\textwidth} >{\centering\arraybackslash}p{0.11\textwidth} >{\centering\arraybackslash}p{0.11\textwidth} >{\centering\arraybackslash}p{0.11\textwidth} >{\raggedright\arraybackslash}X@{}}
        \toprule
        Nastavitev & $k_t$ & $k_s$ & $k_f$ & Opomba \\
        \midrule
        Preiskani prostor & $\{1,2,3\}$ & $\{1,2,3\}$ & $\{1,2,3,5\}$ & $36$ uspešno izvedenih kombinacij \\
        Najvišji makro F1 & $3$ & $1$ & $5$ & makro F1 $0{,}510$, točnost $0{,}522$ \\
        Končna nastavitev & $1$ & $1$ & $3$ & redkejši graf in manjša računska zahtevnost \\
        \bottomrule
    \end{tabularx}
\end{table}


\section{Strojna in programska oprema}
\label{app:strojna-programska-oprema}

Glavni eksperimenti so bili izvedeni na računalniku z operacijskim sistemom Linux
Ubuntu, grafično kartico NVIDIA GeForce RTX 4070 z 12~GB grafičnega pomnilnika,
procesorjem AMD Ryzen 7 5700X, 32~GB sistemskega pomnilnika in pogonom Samsung
SSD 990 PRO NVMe.
Poskusi so bili poganjani v okolju Conda \texttt{gfm}.
Uporabljene so bile različice: Python 3.10.19, PyTorch 2.9.1+cu126, CUDA 12.6, PyG 2.7.0, scikit-learn 1.7.2, LightGBM 4.6.0, NumPy 1.25.2, pandas 2.3.3, Matplotlib 3.10.6 in seaborn 0.13.2.

\section{Naučene uteži povezav}
\label{app:naucene-utezi-povezav}

V modelu \ModelName{HeteroGCN-MLP-w} značilke povezave uporabimo za izračun skalarne uteži povezave.
Naj bo $\mathbf{e}_{ij}^{(r)}$ vektor značilk povezave $(v_i, v_j)$ v relaciji $r$.
Najprej izračunamo surovo oceno:
\begin{equation}
    s_{ij}^{(r)} =
    f_w^{(r)}
    \left(
    \mathbf{e}_{ij}^{(r)}
    \right),
\end{equation}
kjer je $f_w^{(r)}$ večslojni perceptron za ustrezno družino relacij.
Za časovne povezave naprej in nazaj uporabimo isti perceptron $f_w^{\text{čas}}$, za prostorske povezave $f_w^{\text{prostor}}$, za fiksacijske povezave pa $f_w^{\text{fiksacija}}$.
Vsak od teh perceptronov ima obliko
\begin{equation}
    d_e \rightarrow 6 \rightarrow 4 \rightarrow 2 \rightarrow 1,
\end{equation}
kjer je $d_e$ razsežnost vektorja značilk povezave za izbrano relacijo in množico signalov.

Surovo oceno nato omejimo s funkcijo $\tanh$:
\begin{equation}
    \tilde{w}_{ij}^{(r)} =
    \tanh
    \left(
    s_{ij}^{(r)}
    \right).
\end{equation}
Končno utež normaliziramo po vhodnih povezavah ciljne meritve $v_j$ znotraj iste relacije:
\begin{equation}
    w_{ij}^{(r)} =
    \frac{
        \tilde{w}_{ij}^{(r)}
    }{
        \sum_{k : (v_k, v_j) \in E_r}
        \left|
        \tilde{w}_{kj}^{(r)}
        \right|
        + \epsilon
    },
\end{equation}
kjer je $\epsilon=10^{-8}$ majhna konstanta za numerično stabilnost.
Ker normalizacija ohrani predznak uteži, lahko povezava prispevek sosednjega vozlišča tudi zmanjša, ne le poveča.


%%%%%%%%%%%%%%%%%%%%%%%%%%%%%%%%%%%%%%%%
%  LITERATURA
%%%%%%%%%%%%%%%%%%%%%%%%%%%%%%%%%%%%%%%%

\printbibliography[heading=bibintoc]

\end{document}
